\section{Judge Validation}
\label{sec:judgeval}

Replacing expert scores with an LLM judge only helps if the judge is
itself accountable. This section reports four validation experiments on
a frozen 90-question sample, stratified across systems and outcome
labels (\texttt{results/judge\_validation/}), plus a first pass through
the released blinded human-annotation package. Three of the five
results are reassuring; the fifth, agreement with a human reader, is
not---though it does settle which judge model to prefer
(Section~\ref{sec:judgeval:human}).

\subsection{Does the judge discriminate, or merely detect topic?}
\label{sec:judgeval:mismatch}

The sharpest failure mode for an engagement metric is that it measures
topical similarity and calls it engagement. We test this directly:
re-judge every sampled question after swapping in \emph{another
question's} retrieved evidence. A discriminating judge must collapse to
\labelfmt{not\_addressed}.

Engagement falls from 74\% (67/90) on true pairings to 20\% (18/90) on
mismatched ones ($p<10^{-4}$, two-sided Fisher), and the drop is
individually significant for five of six systems. Two further readings
matter. The 20\% residual is a \emph{generosity bound}: on evidence
that cannot possibly bear on the question, this judge still reports
engagement one time in five, so every engagement rate in this paper
should be read against that floor rather than against zero. And the
drop tracks question specificity---cleanest for the most anchored
questions (tension-LLM 12/16 $\to$ 1/16, $p=0.0002$; evidence-graph
8/10 $\to$ 1/10) and weakest for the generic citation-leader template
(11/16 $\to$ 5/16, $p=0.076$, the only non-significant cell). A
question vague enough to accept unrelated evidence is vague enough to
fool the judge, which is independent support for the specificity rubric
of Section~\ref{sec:foresight:saf}.

\subsection{Judge stochasticity and prompt sensitivity}

Temperature 0 is not determinism. Re-running the identical
configuration twice gives outcome agreement 96.7\% and 93.3\%
($\kappa=0.95$, $0.91$) and premise agreement 95.6\% and 94.4\%
($\kappa=0.92$, $0.90$): roughly 3--7 points of label noise, small
relative to the effects the paper reports but not zero, and it should
be assumed present in every rate.

Prompt sensitivity is larger. A semantically equivalent rewrite of the
judge prompt agrees at $\kappa=0.68$ (outcome) and $0.72$ (premise). A
harder variant, replacing the label \emph{names} with neutral codes
(\texttt{L1}--\texttt{L4}, \texttt{P1}--\texttt{P5}) while keeping the
definitions verbatim, drops to $\kappa=0.57$ and $0.55$, and never once
uses the code corresponding to \labelfmt{posed\_but\_open}. Part of the
judge's behaviour therefore rests on the connotations of the label
names, not on their stated definitions. We report this as a real
limitation: label naming is part of the protocol and must be frozen
along with everything else.

\subsection{Cross-model agreement, and what low agreement means here}

Re-judging with \texttt{gpt-4o} and \texttt{gpt-4-turbo} gives low
nominal agreement with \texttt{gpt-4.1}: pairwise Cohen's
$\kappa=0.34$ and $0.20$ on outcome, $0.19$ and $0.05$ on premise
status; three-way Fleiss $\kappa=0.38$ and $0.03$. Taken alone these
numbers say the instrument is unreliable, and we report them
unadorned.

The label distributions complicate that reading. On 90 items
\texttt{gpt-4-turbo} assigns \labelfmt{still\_plausible} 87 times and
never once uses \labelfmt{refuted} or \labelfmt{weakened};
\texttt{gpt-4o} assigns it 78 times. On the outcome dimension both
almost never use \labelfmt{answered} (2 and 4 times of 90, against 25
for \texttt{gpt-4.1}), collapsing a four-way judgement into a
two-way one. Their comparatively high mutual agreement
($\kappa=0.71$ on outcome) is therefore consistent with a shared
conservative default rather than with shared judgement, and $\kappa$ is
in any case depressed when one rater's marginals are near-degenerate.

We flag plainly that this reading is self-serving---``the judges who
disagree with ours are the incompetent ones'' is exactly what a
motivated author would say---and that only human annotation can
arbitrate it. That is the purpose of the released package, and it is
the single most important open item in this paper.

What can be settled without humans is whether the paper's
\emph{conclusions} depend on the judge. Because the stratified sample
equalises label mixes across systems and so erases between-system rate
differences, this requires a second, unstratified draw (40 random
questions from each of four systems); we note the distinction because
computing conclusion robustness on a label-stratified sample is a
mistake that is easy to make and that we made first.

\begin{table}[t]
  \centering
  \small
  \caption{Conclusion robustness under three judges, unstratified
  sample ($n=40$ per system). \checkmark{} = the stated ordering holds;
  $\times$ = it does not. Every $\times$ is a tie at zero, where the
  judge assigns the label to \emph{no} system; no judge reverses any
  conclusion.}
  \label{tab:judgerobust}
  \begin{tabular}{@{}lccc@{}}
    \toprule
    Conclusion & \texttt{gpt-4.1} & \texttt{gpt-4o} & \texttt{gpt-4-turbo} \\
    \midrule
    structure$\to$LLM answers more than LLM-only   & \checkmark & \checkmark & \checkmark \\
    LLM-only engages more than random templates    & \checkmark & \checkmark & \checkmark \\
    structure$\to$LLM refutes more than LLM-only   & \checkmark & \checkmark & $\times$ \\
    structure-only answers more than LLM-only      & \checkmark & $\times$   & $\times$ \\
    \bottomrule
  \end{tabular}
\end{table}

Table~\ref{tab:judgerobust} gives the result. The paper's headline
claim---evidence-structure-first generation resolves more than LLM-only
prompting---holds under all three judges. So does the engagement
ordering. The two conclusions that fail do so in a specific and
benign way: \texttt{gpt-4-turbo} reports 0\% refutation for
\emph{every} system, and both weaker judges report 0\% answered for
both compared systems, so the comparison has no resolution rather than
the opposite sign. Across all twelve judge--conclusion cells, no judge
ever orders the systems the other way.

The practical implication is a requirement, not a reassurance: this
benchmark has a \emph{judge capability floor}. The premise dimension in
particular is unmeasurable with models that default to
\labelfmt{still\_plausible}, so an instance is only reproducible on a
judge that demonstrably uses the full label space. We recommend
reporting the judge's label distribution alongside any submission, and
treating a near-degenerate distribution as a failed run.

\subsection{Human annotation: poor agreement, and an arbitration}
\label{sec:judgeval:human}

The decisive experiment is agreement with a human reader. One annotator,
not a domain specialist, labelled all 90 blinded items from the
abstracts alone, using the released decision procedure and seeing no
system identity and no model label.

\paragraph{The judge is not validated.} Agreement with \texttt{gpt-4.1}
is 34.4\% on outcome ($\kappa=0.10$) and 45.6\% on premise status
($\kappa=0.13$). By any standard reading these are slight, and they do
not establish that the judge measures what a careful reader measures.
Agreement rises with the annotator's confidence but never approaches
adequacy (0\% at confidence 1, 32\% at 2, 41\% at 3, $n=3/53/34$). We
had also hoped the sharpened decision rules would help: the 67 items
labelled after their introduction agree at 35.8\% against 30.4\% for the
first 23, but that difference is not significant ($p=0.80$) and we draw
no conclusion from it.

\begin{table}[t]
  \centering
  \small
  \caption{Human labels against each judge model, $n=90$. Cohen's
  $\kappa$; the human read the same abstracts, blinded to system
  identity and to every model label.}
  \label{tab:humanjudge}
  \begin{tabular}{@{}lcccc@{}}
    \toprule
    & \multicolumn{2}{c}{outcome} & \multicolumn{2}{c}{premise status} \\
    \cmidrule(lr){2-3}\cmidrule(lr){4-5}
    Judge & agreement & $\kappa$ & agreement & $\kappa$ \\
    \midrule
    \texttt{gpt-4.1}      & 34.4\% & \phantom{-}0.10 & 45.6\% & 0.13 \\
    \texttt{gpt-4o}       & 23.3\% & -0.03 & 47.8\% & 0.15 \\
    \texttt{gpt-4-turbo}  & 24.4\% & -0.00 & 42.2\% & 0.04 \\
    \bottomrule
  \end{tabular}
\end{table}

\paragraph{But the cross-model disagreement is now arbitrated.}
Section~\ref{sec:judgeval} reported low agreement between judge models
and offered a reading---that the weaker models collapse to a
conservative default rather than judging---which we flagged as
self-serving and deferred to human evidence. That evidence is now in,
and it supports the reading. On the 42 items where \texttt{gpt-4.1} and
\texttt{gpt-4o} assign different outcomes, the human sides with
\texttt{gpt-4.1} 19 times and with \texttt{gpt-4o} 9 times (with
neither, 14; $p=0.015$ for the 19--9 split). Across all items,
\texttt{gpt-4.1} is the only judge with positive $\kappa$ on outcome
(Table~\ref{tab:humanjudge}), and its label distribution is far closer
to the human's (25/30/12/23 against 33/33/15/9) than either weaker
model's (\texttt{gpt-4o}: 4/47/2/37). The judge capability floor we
proposed on internal evidence is therefore corroborated externally:
models that collapse the label space are worse, not merely different.

\paragraph{The residual disagreement has a direction.} The judge
declares non-engagement far more readily than the human:
\labelfmt{not\_addressed} on 23 of 90 items against the human's 9
($p=0.010$). Of the 23 the judge dismissed, the human found substantive
engagement in 21. Read with the mismatched-evidence control of
Section~\ref{sec:judgeval:mismatch}---where the judge still reported
engagement on 20\% of deliberately unrelated evidence---the instrument
is not simply strict or simply lax: it is \emph{differently} calibrated,
crediting some unrelated evidence while dismissing engagement a human
reader accepts. The confusion is spread across the whole matrix rather
than concentrated in one boundary, which is what one expects when two
raters are applying genuinely different criteria rather than the same
criterion with different thresholds.

\paragraph{What remains open.} Two of the three readings we listed
survive. The judge may be measuring something other than careful human
judgement; or the taxonomy may be underdetermined enough that
independent readers cannot converge on it. What would separate them is
\textbf{human--human agreement}, which requires a second annotator on
the same 90 items and which we do not yet have. The annotator here was
not a domain specialist, which bounds what a single pass can establish
in either direction. Until that number exists, every rate in this paper
should be read as the output of a stable, discriminating, self-consistent
procedure whose correspondence to human judgement is \emph{poor and not
yet explained}. We consider obtaining it, with a domain-expert
annotator, the highest-priority item in this line of work.
